\chapter{A Probability Toolkit}\label{app:A}

This appendix collects the probabilistic facts used in the book, with proofs where they are short and references where they are standard. Readers trained in number theory will find here what they need from probability; readers trained in probability may use it as a list of conventions. Throughout, a law means a Borel probability measure, $\mu_n\Rightarrow\mu$ denotes weak convergence of laws, and $X_n\Rightarrow X$ the same for random variables with those laws.

\section{Weak convergence and the method of moments}

A sequence of laws $\mu_n$ on a metric space $S$ converges weakly to $\mu$ if $\int f\,d\mu_n\to\int f\,d\mu$ for every bounded continuous $f$. By the portmanteau theorem this is equivalent to $\mu_n(A)\to\mu(A)$ for every Borel set $A$ whose boundary is $\mu$ null, and on $\R$ to convergence of distribution functions at continuity points of the limit \cite{billingsley1968}.\index{weak convergence} A law $\mu$ on $\R$ is determined by its characteristic function $\varphi_\mu(t)=\int e^{itx}\,\mu(dx)$, and by L\'evy's continuity theorem $\mu_n\Rightarrow\mu$ if and only if $\varphi_{\mu_n}(t)\to\varphi_\mu(t)$ for every $t\in\R$; more generally, if $\varphi_{\mu_n}$ converges pointwise to a function that is continuous at $0$, that function is the characteristic function of a law $\mu$ and $\mu_n\Rightarrow\mu$ \cite{billingsley1995}.\index{characteristic function}

On a compact interval the method of moments is clean, because polynomials are dense in the continuous functions.

\begin{theorem}[Moments determine compactly supported laws]\label{thm:A:moments}
Let $\mu,\nu$ be laws on a compact interval $I=[a,b]$ with $\int x^k\,d\mu=\int x^k\,d\nu$ for all $k\ge0$. Then $\mu=\nu$. Moreover, if $\mu_n$ are laws on $I$ and $\int x^k\,d\mu_n\to\int x^k\,d\mu$ for every $k$, then $\mu_n\Rightarrow\mu$.
\end{theorem}

\begin{proof}
Let $f\in C(I)$ and $\epsilon>0$. By the Weierstrass approximation theorem there is a polynomial $P$ with $\sup_I|f-P|<\epsilon$. Then
\begin{align*}
\Bigl|\int f\,d\mu-\int f\,d\nu\Bigr|&\le\Bigl|\int (f-P)\,d\mu\Bigr|+\Bigl|\int P\,d\mu-\int P\,d\nu\Bigr|\\
&\quad+\Bigl|\int (P-f)\,d\nu\Bigr|<2\epsilon,
\end{align*}
because the middle term vanishes by linearity. Hence $\int f\,d\mu=\int f\,d\nu$ for all $f\in C(I)$. Continuous functions determine a finite Borel measure on a compact metric space: for a closed set $F$ the functions $f_j(x)=\max(0,1-j\,d(x,F))$ decrease to the indicator of $F$, so by dominated convergence $\mu(F)=\nu(F)$, and closed sets form a $\pi$ system generating the Borel sets, so $\mu=\nu$ by Dynkin's theorem. For the second statement, the same three term estimate gives $\limsup_n|\int f\,d\mu_n-\int f\,d\mu|\le2\epsilon$ for every $\epsilon>0$.
\end{proof}

The theorem applies to the normalised traces $x_p\in[-2,2]$: if the empirical moments converge to those of $\ST$ (the Catalan numbers in even degree, zero in odd degree) then the empirical laws converge weakly to $\ST$. On noncompact spaces moments need not determine laws. Stieltjes \cite{stieltjes1894} showed this for the lognormal law: for $|a|\le1$ the densities $(2\pi)^{-1/2}x^{-1}e^{-(\log x)^2/2}\bigl(1+a\sin(2\pi\log x)\bigr)$ on $(0,\infty)$ all have the moments $e^{k^2/2}$ of the lognormal law. Carleman's condition $\sum_{k\ge1}m_{2k}^{-1/(2k)}=\infty$ on the even moments $m_{2k}$ is a standard sufficient criterion for uniqueness \cite{akhiezer1965}, and if $\mu$ is determined by its moments and $\int x^k\,d\mu_n\to\int x^k\,d\mu$ for every $k$, then $\mu_n\Rightarrow\mu$ \cite{billingsley1995}. Since $\{x^k\}_{k\le m}$ and $\{U_k(x/2)\}_{k\le m}$ span the same space of polynomials, the theorem may equally be stated with Chebyshev moments, which is the form used for $\SU(2)$.\index{method of moments}

\section{Equidistribution, Haar measure and Peter Weyl}

A sequence $(\theta_n)$ in $\R/\Z$ is equidistributed if its empirical laws $N^{-1}\sum_{n\le N}\delta_{\theta_n}$ converge weakly to Lebesgue measure.

\begin{theorem}[Weyl \cite{weyl1916}]\label{thm:A:weyl}
$(\theta_n)$ is equidistributed in $\R/\Z$ if and only if, for every integer $k\ne0$, $N^{-1}\sum_{n\le N}e^{2\pi ik\theta_n}\to0$ as $N\to\infty$.
\end{theorem}

The criterion is the case $\mathrm U(1)\cong\R/\Z$ of Theorem~\ref{thm:04weyl}, in the form for sequences indexed by the integers noted after its proof. Every compact Hausdorff group $G$ carries a unique probability measure $\mu_G$ invariant under left translation, Haar measure; it is also invariant under right translation and inversion. For $\SU(2)$ the trace $x=2\cos\theta$ of a Haar element has the semicircle law $\ST(dx)=\tfrac1{2\pi}\sqrt{4-x^2}\,dx$ on $[-2,2]$, by the Weyl integration formula, which writes the Haar integral of a class function with value $F(\theta)$ on the conjugacy class of eigenvalues $e^{\pm i\theta}$ as $\tfrac2\pi\int_0^\pi F(\theta)\sin^2\theta\,d\theta$.\index{Haar measure}\index{Weyl criterion}

\begin{theorem}[Peter and Weyl \cite{peterweyl1927}]\label{thm:A:pw}
The irreducible unitary representations $\rho$ of a compact group $G$ are finite dimensional, the characters $\chi_\rho$ of the pairwise inequivalent ones form an orthonormal basis of the class functions in $L^2(G,\mu_G)$, and their finite linear combinations are uniformly dense in the continuous class functions.
\end{theorem}

Hence a sequence $(g_n)$ in $G$ has equidistributed conjugacy classes, meaning $N^{-1}\sum_{n\le N}f(g_n)\to\int f\,d\mu_G$ for every continuous class function $f$, if and only if $N^{-1}\sum_{n\le N}\chi_\rho(g_n)\to0$ for every nontrivial irreducible $\rho$. For $\SU(2)$ the irreducible representations are the $\Sym^n$, $n\ge0$, of dimension $n+1$, with characters $U_n(x/2)$, where $U_n(\cos\theta)=\sin((n+1)\theta)/\sin\theta$ is the Chebyshev polynomial of the second kind.

Equidistribution is measured by the Kolmogorov discrepancy. For a continuous distribution function $F$ and the empirical distribution function $F_n$ of $x_1,\dots,x_n$, put $D_n=\sup_t|F_n(t)-F(t)|$. If the $x_i$ are independent with distribution function $F$, then $\sqrt n\,D_n\Rightarrow\mathcal K=\sup_{0\le u\le1}|B(u)|$ for a Brownian bridge $B$, whichever continuous $F$ is used (Kolmogorov \cite{kolmogorov1933}). The Kolmogorov distribution has $\Prob(\mathcal K\le x)=\sum_{k\in\Z}(-1)^ke^{-2k^2x^2}$ for $x>0$, mean $\sqrt{\pi/2}\,\log2=0.8687$, second moment $\pi^2/12$ and standard deviation $0.2603$.\index{Kolmogorov distribution}\index{discrepancy}

\section{Tauberian theorems and infinitely divisible laws}

The Wiener Ikehara theorem \cite{wienerikehara1931,wiener1932}: let $\mu$ be a nonnegative measure on $[0,\infty)$ such that $F(s)=\int e^{-st}\,\mu(dt)$ converges for $\Re s>1$. If $F(s)-c/(s-1)$ extends continuously to the closed half plane $\Re s\ge1$ for some constant $c$, then $e^{-t}\mu([0,t])\to c$ as $t\to\infty$. For a Dirichlet series $F(s)=\sum_nb_nn^{-s}$ with $b_n\ge0$ (take $\mu=\sum_nb_n\delta_{\log n}$) the conclusion is $x^{-1}\sum_{n\le x}b_n\to c$, the form used in the proof of Theorem~\ref{thm:04serre}. The Hardy, Littlewood and Karamata theorem \cite{karamata1931,binghamgoldieteugels1987} states that for a nonnegative measure $\mu$ on $[0,\infty)$ with Laplace transform $\hat\mu(\lambda)=\int e^{-\lambda x}\,\mu(dx)$ finite for $\lambda>0$, a number $\rho\ge0$ and a slowly varying function $\ell$, $\hat\mu(\lambda)\sim\lambda^{-\rho}\ell(1/\lambda)$ as $\lambda\downarrow0$ if and only if $\mu([0,x])\sim x^\rho\ell(x)/\Gamma(1+\rho)$ as $x\to\infty$; the same holds with $\lambda\to\infty$ and $x\to0$.\index{Tauberian theorem}

A law on $\R$ is infinitely divisible if for every $m\ge1$ it is the $m$ fold convolution of a law. Such laws are described by the L\'evy Khintchine theorem \cite{sato1999} (Theorem~\ref{thm:ch09lk}), and for an Euler product $F$ whose logarithm has nonnegative coefficients, $t\mapsto F(\sigma+it)/F(\sigma)$ is a compound Poisson characteristic function at every real $\sigma$ where the series for $\log F$ converges absolutely, as in Theorems~\ref{thm:ch09khinchine} and~\ref{thm:ch09rs}.\index{infinitely divisible law}

\section{Random series}

Kolmogorov's two series theorem: if $Y_k$ are independent real random variables, $\sum_k\E Y_k$ converges and $\sum_k\operatorname{Var}Y_k<\infty$, then $\sum_kY_k$ converges almost surely and in $L^2$. The three series theorem gives the exact criterion: for independent $Y_k$, $\sum_kY_k$ converges almost surely if and only if for some (equivalently every) $c>0$ the three series $\sum_k\Prob(|Y_k|>c)$, $\sum_k\E[Y_k1_{\{|Y_k|\le c\}}]$ and $\sum_k\operatorname{Var}(Y_k1_{\{|Y_k|\le c\}})$ converge \cite{billingsley1995}. For $X=\sum_ka_k\sin\theta_k$ with real $a_k$ and independent phases $\theta_k$ uniform on $[0,2\pi)$, the terms have mean $0$ and variance $a_k^2/2$, so the series converges almost surely and in $L^2$ when $\sum_ka_k^2<\infty$, with $\operatorname{Var}X=\tfrac12\sum_ka_k^2$ and characteristic function $\prod_kJ_0(a_kt)$. Kronecker's theorem links deterministic and random: if $\gamma_1,\dots,\gamma_K$ are linearly independent over $\Q$, then $(\gamma_ky)_{k\le K}$ is equidistributed in $(\R/2\pi\Z)^K$ as $y$ ranges over $[0,Y]$, $Y\to\infty$, which means that the image of normalised Lebesgue measure on $[0,Y]$ under $y\mapsto(\gamma_ky\bmod2\pi)_{k\le K}$ converges weakly to Haar measure. This is the mechanism behind Bohr and Jessen \cite{bohrjessen1930}, Rubinstein and Sarnak \cite{rubinsteinsarnak1994} and the rank bias of Section~\ref{sec:05bias} and Open Problem~\ref{op:rankbias}.\index{random series}

A sequence of real random variables $Y_n$ satisfies a large deviation principle with good rate function $I$ if $I:\R\to[0,\infty]$ is lower semicontinuous with compact sublevel sets, $\limsup_nn^{-1}\log\Prob(Y_n\in C)\le-\inf_CI$ for every closed $C$, and $\liminf_nn^{-1}\log\Prob(Y_n\in O)\ge-\inf_OI$ for every open $O$. By the G\"artner Ellis theorem \cite[Theorem 2.3.6]{dembozeitouni1998}, if $\Lambda(\theta)=\lim_nn^{-1}\log\E e^{n\theta Y_n}$ exists in $(-\infty,\infty]$ for every real $\theta$, if $0$ is an interior point of $\{\Lambda<\infty\}$, and if $\Lambda$ is lower semicontinuous, differentiable on the interior of its domain and steep, meaning $|\Lambda'(\theta)|\to\infty$ as $\theta$ approaches a finite endpoint of the domain from inside, then $Y_n$ satisfies the large deviation principle with the good rate function $\Lambda^*(y)=\sup_\theta(\theta y-\Lambda(\theta))$.\index{large deviations}

\section{Generalized gamma convolutions and P\'olya frequency functions}

By Bernstein's theorem \cite{bernstein1929,widder1941} a function $f$ on $(0,\infty)$ is completely monotone, $(-1)^nf^{(n)}\ge0$ for all $n\ge0$, if and only if it is the Laplace transform $f(x)=\int_{[0,\infty)}e^{-xt}\,\tau(dt)$ of a nonnegative measure $\tau$, which is then unique.\index{Bernstein's theorem} A law on $[0,\infty)$ is a generalized gamma convolution (GGC) if it is a weak limit of finite convolutions of gamma laws \cite{thorin1977,bondesson1992}; equivalently
\[
\E e^{-sX}=\exp\Bigl(-as-\int_{(0,\infty)}\log\Bigl(1+\frac{s}{t}\Bigr)\,\tau(dt)\Bigr),\qquad s\ge0,
\]
with $a\ge0$ and a positive measure $\tau$ on $(0,\infty)$, the Thorin measure, such that $\int\log(1+1/t)\,\tau(dt)<\infty$ (Definition~\ref{def:ch10ggc}). By Frullani's integral a GGC is infinitely divisible with drift $a$ and L\'evy density $x^{-1}\int e^{-tx}\,\tau(dt)$, and conversely every infinitely divisible law on $[0,\infty)$ of this form is a GGC; the sum $\sum_kE_k/c_k$ of independent standard exponentials, with $c_k>0$ and $\sum_kc_k^{-1}<\infty$, has Thorin measure $\sum_k\delta_{c_k}$; and a GGC Laplace transform extends holomorphically to $\C\setminus(-\infty,0]$ without zeros (Lemma~\ref{lem:ch10nonvanish}), which drives the Thorin reformulation of Part~IV.\index{generalized gamma convolution}

A function $f$ on $\R$ is a P\'olya frequency function if $\det[f(x_i-y_j)]_{i,j\le m}\ge0$ for every $m\ge1$ and all $x_1<\dots<x_m$, $y_1<\dots<y_m$. Schoenberg's theorem \cite{schoenberg1951,karlin1968} (Theorem~\ref{thm:ch10schoenberg}) characterises the integrable ones, not almost everywhere zero, by their bilateral Laplace transforms $1/\Psi$, with $\Psi$ of Laguerre P\'olya type; in its notation, $f/\!\int f$ is the density of $\delta+\sqrt{2\gamma}\,Z+\sum_k\delta_k(E_k-1)$, with $Z$ standard normal, the $E_k$ standard exponential, and $Z,E_1,E_2,\dots$ independent.\index{P\'olya frequency function}

\section{Extreme values}

The maximum $M_n$ of $n$ independent random variables with distribution function $F$ has distribution function $F^n$. By the Fisher Tippett Gnedenko theorem \cite{fishertippett1928,gnedenko1943,leadbetter1983}, if $(M_n-b_n)/a_n$ converges in law to a nondegenerate limit for some constants $a_n>0$ and $b_n$, then the limit is, after a change of location and scale, a Gumbel law $\exp(-e^{-x})$, a Fr\'echet law $\exp(-x^{-\alpha})$ on $x>0$, or a Weibull law $\exp(-(-x)^\alpha)$ on $x<0$, with $\alpha>0$.\index{extreme value theory} By the R\'enyi representation \cite{renyi1953} the maximum of $n$ independent standard exponentials has the law of $\sum_{j\le n}E_j'/j$, with $E_j'$ independent standard exponentials, a GGC with Thorin measure $\sum_{j\le n}\delta_j$. The infinite version used in the book is
\[
\Prob\Bigl(\max_{n\ge1}\frac{E_n}{cn}\le y\Bigr)=\prod_{n\ge1}\bigl(1-e^{-cny}\bigr),\qquad y>0,
\]
for $c>0$; the maximum is almost surely finite because $\sum_ne^{-cny}<\infty$. With $c=2\pi$ the product is $e^{\pi y/12}\eta(iy)$, the starting point of Chapter~\ref{ch:11} (Proposition~\ref{prop:ch11M} and \eqref{eq:ch11DeltaF}).

\section{Markov chains on graphs}

Let $P$ be the transition matrix of an irreducible Markov chain on a finite set $V$ of size $n$ with stationary law $\pi$, so $\pi(x)>0$ for all $x$. The chain is reversible if $\pi(x)P(x,y)=\pi(y)P(y,x)$. Simple random walk on a connected $d$ regular graph, whose adjacency matrix $A$ is a symmetric matrix of nonnegative integers with all row sums $d$ (loops and multiple edges are allowed), has $P=A/d$ and is reversible with uniform $\pi$.\index{Markov chain}

\paragraph{Spectral decomposition.} If $P$ is reversible, then $D^{1/2}PD^{-1/2}$ with $D=\mathrm{diag}(\pi)$ is symmetric, so $P$ has real eigenvalues $1=\mu_1>\mu_2\ge\dots\ge\mu_n\ge-1$ and eigenvectors $\phi_1=1,\phi_2,\dots,\phi_n$ orthonormal in $L^2(\pi)$. Then
\[
P^t(x,y)=\pi(y)\sum_{j=1}^n\mu_j^t\,\phi_j(x)\phi_j(y).
\]
The strict inequality $\mu_2<1$ is irreducibility, and $\mu_n=-1$ if and only if the chain has period $2$ (for a connected graph, if and only if it is bipartite). Averaged return probabilities $n^{-1}\tr P^t=n^{-1}\sum_j\mu_j^t$ are the moments of the spectral measure.

\paragraph{Hitting times.} Let $\tau_y=\min\{t\ge0:X_t=y\}$ and $h(x,y)=\E_x\tau_y$. The function $h(\cdot,y)$ is the unique solution of the system $h(y,y)=0$ and, for $x\ne y$, $h(x,y)=1+\sum_zP(x,z)h(z,y)$. With $\tau_y^+=\min\{t\ge1:X_t=y\}$, Kac's formula gives the mean return time $\E_y\tau_y^+=1/\pi(y)$ \cite{aldousfill}.

\paragraph{Mixing and the spectral gap.} The spectral decomposition also controls convergence to equilibrium. With $\mu_*=\max(\mu_2,|\mu_n|)$, the Cauchy Schwarz inequality gives
\[
\Bigl|\frac{P^t(x,y)}{\pi(y)}-1\Bigr|\le\mu_*^t\Bigl(\sum_{j\ge2}\phi_j(x)^2\Bigr)^{1/2}\Bigl(\sum_{j\ge2}\phi_j(y)^2\Bigr)^{1/2}\le\frac{\mu_*^t}{\sqrt{\pi(x)\pi(y)}},
\]
since $\sum_j\phi_j(x)^2=1/\pi(x)$ by orthonormality. For a $d$ regular graph this gives $|P^t(x,y)-1/n|\le\mu_*^t$, which is empty for a bipartite graph, where $\mu_*=1$. A connected Ramanujan graph that is not bipartite, so that with $q=d-1$ every adjacency eigenvalue $\lambda$ with $|\lambda|\ne d$ has $|\lambda|\le2\sqrt q$ and $-d$ is not an eigenvalue, has $\mu_*\le\rho_q=2\sqrt q/(q+1)$, the spectral radius of simple random walk on the $d$ regular tree \cite{kesten1959}; hence $|P^t(x,y)-1/n|\le\rho_q^t$.

\paragraph{Kemeny's constant.} Kemeny and Snell \cite{kemenysnell1960} observed that the quantity $K=\sum_y\pi(y)h(x,y)$ does not depend on the starting state $x$. This is Kemeny's constant: the expected time to hit a target drawn from the stationary law. In our convention $h(y,y)=0$; putting the mean return time $1/\pi(y)$ on the diagonal instead would add $1$ to $K$.\index{Kemeny constant} Let $\Pi$ be the matrix with every row equal to $\pi$, and let $Z=\sum_{t\ge0}(P^t-\Pi)$ be the fundamental matrix. By the spectral decomposition $P^t(x,y)-\pi(y)=\pi(y)\sum_{j\ge2}\mu_j^t\phi_j(x)\phi_j(y)$, so if $\mu_n>-1$ the series converges and summing the geometric series gives
\[
Z(x,y)=\pi(y)\sum_{j\ge2}\frac{\phi_j(x)\phi_j(y)}{1-\mu_j}.
\]
If $\mu_n=-1$ the chain is periodic, the series is understood in the Abel sense, $Z=\lim_{r\uparrow1}\sum_tr^t(P^t-\Pi)$, and the displayed formula remains valid because $(1-r\mu_j)^{-1}\to(1-\mu_j)^{-1}$ for every $j\ge2$. In all cases we take the displayed formula as the definition of $Z$.

\begin{theorem}\label{thm:A:kemeny}
Let $P$ be irreducible and reversible with eigenvalues $1=\mu_1>\mu_2\ge\dots\ge\mu_n\ge-1$. Then for all $x,y$
\[
h(x,y)=\frac{Z(y,y)-Z(x,y)}{\pi(y)},
\]
and for every $x$
\[
K=\sum_y\pi(y)\,h(x,y)=\sum_{j=2}^n\frac{1}{1-\mu_j}.
\]
In particular, for simple random walk on a connected $d$ regular graph with adjacency eigenvalues $d=\lambda_1>\lambda_2\ge\dots\ge\lambda_n$, $h(x,y)=n\bigl(Z(y,y)-Z(x,y)\bigr)$ and $K=\sum_{j\ge2}(1-\lambda_j/d)^{-1}$.
\end{theorem}

\begin{proof}
The matrix $(\sqrt{\pi(x)}\,\phi_j(x))_{x,j}$ is orthogonal, so its rows are orthonormal as well:
\[
\sum_{j=1}^n\pi(y)\phi_j(x)\phi_j(y)=1_{x=y}.
\]
Since $\phi_1=1$, this gives
\[
\sum_{j\ge2}\pi(y)\phi_j(x)\phi_j(y)=1_{x=y}-\pi(y).
\]
Using $P\phi_j=\mu_j\phi_j$ and $\mu_j/(1-\mu_j)=1/(1-\mu_j)-1$ we find
\[
(PZ)(x,y)=\pi(y)\sum_{j\ge2}\frac{\mu_j\phi_j(x)\phi_j(y)}{1-\mu_j}=Z(x,y)-1_{x=y}+\pi(y).
\]
Fix $y$ and define $g(x)=(Z(y,y)-Z(x,y))/\pi(y)$. Then $g(y)=0$, and for $x\ne y$,
\begin{align*}
\sum_zP(x,z)g(z)&=\frac{Z(y,y)-(PZ)(x,y)}{\pi(y)}\\
&=\frac{Z(y,y)-Z(x,y)-\pi(y)}{\pi(y)}=g(x)-1.
\end{align*}
So $g$ solves the hitting time equations, and by uniqueness $g=h(\cdot,y)$. (Uniqueness: the difference $u$ of two solutions vanishes at $y$ and is harmonic off $y$. If $\max u>0$, it is attained at some $x\ne y$, harmonicity makes every $z$ with $P(x,z)>0$ a maximum point as well, and irreducibility carries the maximum along a path to $y$, a contradiction; the same argument applies to $-u$.) Now
\[
\sum_y\pi(y)h(x,y)=\sum_yZ(y,y)-\sum_yZ(x,y).
\]
The second sum vanishes, because $\sum_y\pi(y)\phi_j(y)=\langle\phi_j,\phi_1\rangle_\pi=0$ for $j\ge2$. The first is the trace of $Z$, and $\sum_y\pi(y)\phi_j(y)^2=1$ by orthonormality, so
\[
\tr Z=\sum_{j\ge2}\frac{1}{1-\mu_j}\sum_y\pi(y)\phi_j(y)^2=\sum_{j\ge2}\frac{1}{1-\mu_j}.
\]
This is independent of $x$, and the regular graph case follows with $\pi(y)=1/n$ and $\mu_j=\lambda_j/d$.
\end{proof}

\section{Divergence, Bayes factors and consistency}

The Kullback Leibler divergence \cite{kullbackleibler1951} $\mathrm{KL}(\mu\|\nu)=\int\log(d\mu/d\nu)\,d\mu$, set equal to $+\infty$ unless $\mu\ll\nu$, is nonnegative, zero only for $\mu=\nu$, not symmetric, and jointly convex in $(\mu,\nu)$.\index{Kullback Leibler divergence} For independent observations $X_1,X_2,\dots$ from $f_0$, the Bayes factor of $f_1$ against $f_0$ \cite{kassraftery1995} is $B_n=\prod_{i\le n}f_1(X_i)/f_0(X_i)$, and by the strong law of large numbers
\[
\frac1n\log B_n\to-\mathrm{KL}(f_0\|f_1)\qquad\text{almost surely},
\]
with limit $-\infty$ if the divergence is infinite. This is the principle behind the rates of Chapter~\ref{ch:12}, where equidistribution replaces the law of large numbers.\index{Bayes factor} Wald's sequential probability ratio test~\cite{wald1947}, whose stopping rule Chapter~\ref{ch:12} states with the target probabilities $\alpha$ of rejecting $f_0$ when it holds and $\beta$ of accepting it when $f_1$ holds, needs for small $\alpha$ and $\beta$, ignoring overshoot, about $\log(1/\beta)/\mathrm{KL}(f_0\|f_1)$ observations under $f_0$ and $\log(1/\alpha)/\mathrm{KL}(f_1\|f_0)$ under $f_1$, and is optimal in expected sample size among tests with no larger error probabilities~\cite{waldwolfowitz1948}. Doob's theorem \cite{doob1949} gives posterior consistency at almost every parameter value under the prior, for a parametric model with a measurable and identifiable parametrisation, and the maximum likelihood estimator is consistent under Wald's conditions \cite{wald1949}.\index{posterior consistency} If the true density is one of finitely many candidate densities $f_1,\dots,f_m$, say $f_t$, with $0<\mathrm{KL}(f_t\|f_i)$ for $i\ne t$, then by the law of large numbers above the posterior probability of each wrong model $i$ is $\exp(-n(\mathrm{KL}(f_t\|f_i)+o(1)))$ almost surely. If the truth $f_0$ lies outside the list, the posterior concentrates on the candidates that minimise $\mathrm{KL}(f_0\|f_i)$ \cite{berk1966}, and the exponential rates are differences of divergences.

\section{Notation}\label{sec:Anotation}

The table lists the symbols that are used in more than one chapter or are central to the argument of a chapter, with the place where each is defined, while symbols that serve a single argument are defined where they occur.

{\footnotesize\setlength{\tabcolsep}{3pt}
\begin{longtable}{p{0.8in}p{2.65in}p{0.95in}}
\toprule
Symbol & Meaning & Where defined\\
\midrule
\endfirsthead
\toprule
Symbol & Meaning & Where defined\\
\midrule
\endhead
\multicolumn{3}{l}{\textit{Number theory and $L$ functions}}\\*
\raggedright $a_p$ & \raggedright Trace of Frobenius at a good prime $p$, $a_p=p+1-\#E(\mathbb F_p)$ for an elliptic curve $E$; the Hecke eigenvalue $a(p)$ of a newform & \raggedright Chapter~\ref{ch:01}; Section~\ref{sec:ch3artin}\tabularnewline[1pt]
\raggedright $x_p$, $\theta_p$ & \raggedright Normalised trace $x_p=a_p/p^{(k-1)/2}\in[-2,2]$ ($a_p/\sqrt p$ in weight $2$), $x_p=2\cos\theta_p$ with $\theta_p\in[0,\pi]$, and $x_p=\tr A_p(\pi)$ in the unitary normalisation; $\theta_p$ is also the Frobenius class in $K^\natural$ and an independent Sato Tate angle in Chapters~\ref{ch:09}, \ref{ch:09b}, \ref{ch:10c} & \raggedright Chapter~\ref{ch:01}; Theorem~\ref{thm:01satotate}; Theorem~\ref{thm:04weyl}\tabularnewline[1pt]
\raggedright $\pi$ & \raggedright Cuspidal automorphic representation of $\GL_n(\A_\Q)$ with unitary central character, contragredient $\tilde\pi$; $\pi_f$ is the representation of a newform $f$ & \raggedright Chapter~\ref{ch:02}\tabularnewline[1pt]
\raggedright $\pi(x)$ & \raggedright Number of primes $p\le x$ & \raggedright Chapter~\ref{ch:01}\tabularnewline[1pt]
\raggedright $Q(x)$, $Q_\alpha(x)$, $D(s)$ & \raggedright $Q(x)=\sum_{n\le x}\mu(n)^2$, the number of squarefree $n\le x$, and $D(s)=\sum_n\mu(n)^2n^{-s}=\zeta(s)/\zeta(2s)$; $Q_\alpha(x)$ is the sum of $\mu^2(\mathfrak a)N(\mathfrak a)^{-\alpha}$ over the ideals $\mathfrak a$ of norm at most $x$ of a number field; Chapter~\ref{ch:squarefree} writes $\Delta(x)=Q(x)-6x/\pi^2$, which is not Ramanujan's $\Delta$ & \raggedright Chapter~\ref{ch:squarefree}\tabularnewline[1pt]
\raggedright $A_p(\pi)$, $\alpha_{j,p}$ & \raggedright Satake parameter of $\pi$ at an unramified prime $p$, a semisimple class in $\GL_n(\C)$ with eigenvalues $\alpha_{j,p}$ & \raggedright Chapter~\ref{ch:02}; Theorem~\ref{thm:ch3satake}\tabularnewline[1pt]
\raggedright $L(s,\pi)$ & \raggedright Standard $L$ function $\prod_p\det(I-A_p(\pi)p^{-s})^{-1}$ in the unitary (analytic) normalisation, critical line $\Re s=\frac12$; the default for $\pi$ and for $\Sym^nf$ & \raggedright Chapter~\ref{ch:01}; Chapter~\ref{ch:02}\tabularnewline[1pt]
\raggedright $L(s,f)$, $L(s,\Delta)$ & \raggedright Classical normalisation $\sum_na(n)n^{-s}$ of a newform of weight $k$, centre $s=\frac k2$, so $L(s,\pi_f)=L(s+\frac{k-1}2,f)$; for $\Delta$ it is $\sum_n\tau(n)n^{-s}$, centre $s=6$ & \raggedright Theorem~\ref{thm:02fe}; Chapter~\ref{ch:10}\tabularnewline[1pt]
\raggedright $L(s,\Delta)$ & \raggedright Unitary normalisation $\sum_n\tau(n)n^{-11/2-s}$, centre $s=\frac12$; it is the classical $L(s+\frac{11}2,\Delta)$ & \raggedright Chapter~\ref{ch:09}; Chapter~\ref{ch:10c}\tabularnewline[1pt]
\raggedright $L(s,E)$ & \raggedright $L$ function of an elliptic curve $E/\Q$, $\prod_p(1-a_pp^{-s}+p^{1-2s})^{-1}$ at good $p$, in the classical normalisation with centre $s=1$ & \raggedright Section~\ref{sec:02level}; Theorem~\ref{thm:02modularity}\tabularnewline[1pt]
\raggedright $E_1$, $E_2$ & \raggedright The curves $y^2=x^3+x+1$ (conductor $496$, no complex multiplication) and $y^2=x^3-x$ (complex multiplication by $\Z[i]$, conductor $32$); in Chapter~\ref{ch:hit}, $E_1(v)=\int_v^\infty e^{-x}x^{-1}dx$ is the exponential integral & \raggedright Chapter~\ref{ch:01}; Figure~\ref{fig:satotate}; Theorem~\ref{thm:hitrate}\tabularnewline[1pt]
\raggedright $\Lambda(s,\pi)$ & \raggedright Completed $L$ function $\gamma(s,\pi)L(s,\pi)=\varepsilon(\pi)\Lambda(1-s,\tilde\pi)$ with $\gamma(s,\pi)=N_\pi^{s/2}\prod_j\Gamma_\R(s+\mu_j)$ and $\Gamma_\R(s)=\pi^{-s/2}\Gamma(\frac s2)$, $\Gamma_\C(s)=\Gamma_\R(s)\Gamma_\R(s+1)=2(2\pi)^{-s}\Gamma(s)$; for an elliptic curve, $\Lambda(s,E)=N_E^{s/2}(2\pi)^{-s}\Gamma(s)L(s,E)$, centre $1$ (Chapters~\ref{ch:02} and \ref{ch:flt}) & \raggedright Section~\ref{sec:10bcompletion}; Chapter~\ref{ch:02}; Corollary~\ref{cor:recipconverse}\tabularnewline[1pt]
\raggedright $\Lambda(s,\Delta)$ & \raggedright Completed function $(2\pi)^{-s}\Gamma(s)L(s,\Delta)=\Lambda(12-s,\Delta)$, centre $6$ (classical); $(2\pi)^{-s-11/2}\Gamma(s+\frac{11}2)L(s,\Delta)$, centre $\frac12$ (unitary; half of $\Lambda(s,\pi)$ of the row above, by Legendre duplication) & \raggedright Theorem~\ref{thm:02fe}; Chapter~\ref{ch:10}; Section~\ref{sec:10ccombs}\tabularnewline[1pt]
\raggedright $\beta(s)$, $\chi_{-4}$, $\chi_{12}$ & \raggedright Dirichlet's $\beta(s)=L(s,\chi_{-4})$, completed as $\Lambda(s)=(4/\pi)^{s/2}\Gamma(\frac{s+1}2)\beta(s)$, and $\Lambda_{12}(s)=(12/\pi)^{s/2}\Gamma(\frac s2)L(s,\chi_{12})$, with P\'olya kernels built from $\eta^3$ and $\eta$ & \raggedright Section~\ref{beta:sec:primal}; Section~\ref{beta:sec:eta}\tabularnewline[1pt]
\raggedright $\varepsilon(\pi)$, $\varepsilon_f$, $\varepsilon(\chi)$, $w(E)$ & \raggedright Root number, $|\varepsilon(\pi)|=1$; classically $\Lambda(s,f)=\varepsilon_f\Lambda(k-s,f)$, $\varepsilon_f=\pm i^k$; $\varepsilon(\chi)$ for a Dirichlet character, $w(E)=\varepsilon_f=\pm1$ for an elliptic curve $E$ & \raggedright Theorem~\ref{thm:02fe}; Section~\ref{sec:10bcompletion}; Chapter~\ref{ch:06c}\tabularnewline[1pt]
\raggedright $g(\chi)$, $g(m)$ & \raggedright Gauss sum $g(\chi)=\sum_{a\bmod N}\chi(a)e^{2\pi ia/N}$ of a Dirichlet character modulo $N$; for primitive $\chi$ with $\chi(-1)=(-1)^\delta$, $\varepsilon(\chi)=g(\chi)/(i^\delta\sqrt N)$; quadratic Gauss sum $g(m)=\sum_{n\bmod m}e^{2\pi in^2/m}$ for odd $m$ & \raggedright Chapter~\ref{ch:02}; Corollary~\ref{cor:recipgausssign}\tabularnewline[1pt]
\raggedright $\Delta$, $\tau(n)$, $q$ & \raggedright $\Delta(\tau)=q\prod_{n\ge1}(1-q^n)^{24}=\sum_n\tau(n)q^n$, the weight $12$ cusp form ($\tau$ in the upper half plane, nome $q=e^{2\pi i\tau}$, Ramanujan function $\tau(n)$); $\tau(p)/p^{11/2}=2\cos\theta_p$ & \raggedright Section~\ref{sec:02forms}; Example~\ref{ex:02tau}\tabularnewline[1pt]
\raggedright $G_\Q$, $\Frob_p$ & \raggedright Absolute Galois group $\Gal(\overline{\Q}/\Q)$, and the Frobenius class at an unramified prime $p$ & \raggedright Section~\ref{sec:ch3cheb}\tabularnewline[1pt]
\raggedright $N$ & \raggedright Conductor $N_\pi$ or $N_E$, or level of a newform in $S_k(\Gamma_0(N))$ & \raggedright Section~\ref{sec:02level}\tabularnewline[1pt]
\raggedright $S_k(\Gamma_0(N))$, $T_n$ & \raggedright Cusp forms of weight $k$ and level $N$, and the Hecke operators $T_n$, with $T_p$ acting on a newform by $a(p)$ & \raggedright Proposition~\ref{prop:02hecke}; Section~\ref{sec:02level}\tabularnewline[1pt]
\raggedright $\lambda_f(p)$ & \raggedright Unitary Hecke eigenvalue $a_f(p)/p^{(k-1)/2}=x_p\in[-2,2]$ of a newform $f$; $\lambda(p)$ for $\Delta$ and $\lambda_\phi(p)$ for a Maass form $\phi$ & \raggedright Section~\ref{sec:ch6ils}; Open Problem~\ref{op:rp}\tabularnewline[1pt]
\raggedright $\Lambda(n)$, $\Lambda_E(n)$, $\Lambda_\pi(n)$, $\Lambda_\rho(n)$ & \raggedright The von Mangoldt function; coefficients of $-L'/L$: $\Lambda_E(p^k)=p^{k/2}\,2\cos(k\theta_p)\log p$ (classical, centre $1$) and $-\frac{L'}L(s,\pi)=\sum_n\Lambda_\pi(n)n^{-s}$ (unitary), and for an Artin representation $\rho$, $\Lambda_\rho(p^m)=\tr\rho(\Frob_p^m)\log p$ at unramified $p$ & \raggedright Chapter~\ref{ch:04}; Section~\ref{sec:05bias}; Chapter~\ref{ch:10d}; Definition~\ref{def:recipheat}\tabularnewline[1pt]
\raggedright $\pi_{E,t}$, $C_{E,t}$ & \raggedright Lang Trotter count $\pi_{E,t}(x)={\#\{p\le x:p\nmid N,\ a_p=t\}}\sim C_{E,t}\sqrt x/\log x$ & \raggedright Chapter~\ref{ch:05}; Open Problem~\ref{op:lt}\tabularnewline[1pt]
\raggedright $S(x)$, $r$ & \raggedright $S(x)=\sum_{p\le x,\,p\nmid N}a_p\log p/p\approx(\frac12-r)\log x$ under the Riemann hypothesis for $L(s,E)$ (rank bias), where $r=\operatorname{ord}_{s=1}L(s,E)$ is the analytic rank; in Chapter~\ref{ch:squarefree}, $S(x)=\sum_n\mu(n)^2e^{-n/x}$ & \raggedright Section~\ref{sec:05bias}; Open Problem~\ref{op:rankbias}\tabularnewline[1pt]
\raggedright $\gamma_k$, $\gamma_{\mathrm E}$, $\gamma_K$ & \raggedright Ordinates of zeros: $\frac12\pm i\gamma_k$ for $\Lambda(s,\pi)$ (unitary), $6\pm i\gamma_k$ for $\Lambda(s,\Delta)$ (classical), all real exactly under the Riemann hypothesis; $\gamma_{\mathrm E}$ is Euler's constant; $\gamma_K$ is Ihara's Euler Kronecker constant of a number field $K$ & \raggedright Equation \eqref{eq:10batoms}; Theorem~\ref{thm:ch10hadamard}; Section~\ref{sec:05bias}; Theorem~\ref{thm:recipdrift}\tabularnewline[1pt]
\raggedright $\kappa$, $\kappa_K$, $c_K$, $\kappa_C$ & \raggedright In Chapter~\ref{ch:squarefree}, $\sum_{n\le x}\mu(n)^2/n=\frac6{\pi^2}(\log x+\kappa)+O(x^{-1/2})$ with $\kappa=24\log A-\gamma_{\mathrm E}-2\log2\pi$, $A$ the Glaisher Kinkelin constant; over a number field $K$, $c_K=\operatorname{Res}_{s=1}\zeta_K(s)$ and the constant becomes $\kappa_K=\gamma_K-2\zeta_K'(2)/\zeta_K(2)$; in Chapters~\ref{ch:recip} and \ref{ch:pd}, $\kappa_C$ is the Euler Kronecker drift of a conjugacy class $C$; other chapters use $\kappa$ as a local letter & \raggedright Proposition~\ref{prop:sf:classical}; Theorem~\ref{thm:sf:const}; Theorem~\ref{thm:recipdrift}\tabularnewline[1pt]
\raggedright $\delta(C)$, $b_C$, $c_C$, $h_C$, $K_C$ & \raggedright In Chapter~\ref{ch:pd}, for a conjugacy class $C$ of $G=\Gal(K/\Q)$: its Haar measure $\delta(C)=|C|/|G|$ (written $\mu(C)$ in Chapter~\ref{ch:recip}); the Mertens constants $b_C$ and $c_C$, the limits of $\sum_{p\le t,\,p\in C}1/p-\delta(C)\log\log t$ and of $\sum_{p\le t,\,p\in C}\log p/p-\delta(C)\log t$ over the unramified primes; $h_C(s)=P_C(s)-\delta(C)\log\frac1{s-1}$, as in the proof of Theorem~\ref{thm:recipdrift}; and $K_C=e^{-h_C(1)}/\Gamma(1-\delta(C))$, the constant of the proportion $K_C(\log x)^{-\delta(C)}$ of the integers $n\le x$ with no prime factor in $C$ & \raggedright Theorems~\ref{thm:pdmertens}, \ref{thm:pdmertensL} and \ref{thm:pdrate}\tabularnewline[1pt]
\midrule
\multicolumn{3}{l}{\textit{Groups and Haar measure}}\\*
\raggedright $h$, $h^\natural$, $\chi_\rho$ & \raggedright Normalised Haar measure of a compact group $K$ (written $dh$ in Chapter~\ref{ch:03}), its image $h^\natural$ on the space $K^\natural$ of conjugacy classes, and the character $\chi_\rho=\tr\rho$ & \raggedright Chapter~\ref{ch:04}; Theorem~\ref{thm:04weyl}\tabularnewline[1pt]
\raggedright $\mu_G$ & \raggedright Haar probability measure of a compact group $G$; for $G=\mathrm{ST}(A)\subset\USp(2g)$ also its image under the trace, a law on $[-2g,2g]$ & \raggedright Theorem~\ref{thm:A:pw}; Chapter~\ref{ch:12}\tabularnewline[1pt]
\raggedright $\bar\rho_{E,p}$, $G_{E,p}$, $\chi_p$ & \raggedright Mod $p$ Galois representation of an elliptic curve $E$, its image $G_{E,p}\subset\GL_2(\mathbb F_p)$ (written $G(\ell)$ in Chapter~\ref{ch:05}), and the mod $p$ cyclotomic character $\chi_p=\det\bar\rho_{E,p}$ & \raggedright Section~\ref{sec:fltproof}; Definition~\ref{def:flteis}; Proposition~\ref{prop:flteisgamma}\tabularnewline[1pt]
\raggedright $a_\rho$, $\Prob_\sigma$, $d(\rho,\rho')$ & \raggedright Frobenius trace variable $a_\rho(p^m)=\tr\rho(\Frob_p^m)$ on the prime powers $p^m$, $p\notin S$, under $\Prob_\sigma(\{p^m\})\propto p^{-m\sigma}/m$, and $a_\pi(p^m)$, the power sum of the Satake parameters of a cuspidal $\pi$; the Frobenius distance $d(\rho,\rho')=\lim_{\sigma\downarrow1}(\E_\sigma|a_\rho-a_{\rho'}|^2)^{1/2}$ & \raggedright Section~\ref{sec:reciplaw}; Theorem~\ref{thm:recipcheb}\tabularnewline[1pt]
\raggedright $\ST$ & \raggedright Sato Tate (semicircle) law $\frac1{2\pi}\sqrt{4-x^2}\,dx$ on $[-2,2]$: the law of $x=2\cos\theta$ under Haar measure on $\SU(2)$ & \raggedright Chapter~\ref{ch:01}; Proposition~\ref{prop:04catalan}\tabularnewline[1pt]
\raggedright $U_n$, $\Sym^n$ & \raggedright Chebyshev polynomial of the second kind, $U_n(\cos\theta)=\sin((n+1)\theta)/\sin\theta$; $\tr\Sym^n=U_n(x/2)$ is the character of the $n$th symmetric power of $\SU(2)$, orthonormal in $L^2(\ST)$ & \raggedright Proposition~\ref{prop:04catalan}; Theorem~\ref{thm:A:pw}\tabularnewline[1pt]
\raggedright $\mu_{\mathrm{CM}}$ & \raggedright Complex multiplication law $\frac12\delta_0+\frac12\,dx/(\pi\sqrt{4-x^2})$ on $[-2,2]$; $\int U_n(x/2)\,d\mu_{\mathrm{CM}}$ is $1$ if $4\mid n$ and $0$ otherwise & \raggedright Section~\ref{sec:04cm}; Proposition~\ref{prop:04cm}\tabularnewline[1pt]
\raggedright $\mathrm{ST}(A)$ & \raggedright Sato Tate group of an abelian variety $A/\Q$ of dimension $g$, a compact subgroup of $\USp(2g)$; $\SU(2)$ or $N(\mathrm U(1))$ for an elliptic curve & \raggedright Definition~\ref{def:ch3stgroup}; Conjecture~\ref{conj:ch3gst}\tabularnewline[1pt]
\raggedright $s_p$, $W_\rho(n)$ & \raggedright Normalised Frobenius class of $A$ at $p$ in $G=\mathrm{ST}(A)$, $x_p=\tr s_p\in[-2g,2g]$; Weyl average $W_\rho(n)=n^{-1}\sum_{k\le n}\chi_\rho(s_{p_k})$, with limit $\int\chi_\rho\,d\mu_G$ (the Weyl fingerprint); written $W_n(x)$ in Chapter~\ref{ch:04} and $W_m$ in Chapter~\ref{ch:05}, and unrelated to the P\'olya variable $W$ & \raggedright Section~\ref{sec:ch3stgroup}; Chapter~\ref{ch:12}\tabularnewline[1pt]
\midrule
\multicolumn{3}{l}{\textit{Random matrices, families and random models}}\\*
\raggedright $Z_N(\theta)$, $\Lambda_N$ & \raggedright Characteristic polynomial $\det(I-U_Ne^{-i\theta})$ of a Haar unitary matrix $U_N$ of size $N$, and its value $\Lambda_N=Z_N(0)=\det(I-U_N)$ at $\theta=0$ & \raggedright Section~\ref{sec:ch6ksnaith}; Chapter~\ref{ch:09b}, Proposition~\ref{prop:09bks}\tabularnewline[1pt]
\raggedright $N$ & \raggedright Size of a random matrix in $\mathrm U(N)$, $\USp(2N)$ or $\mathrm{SO}(2N)$, matched to arithmetic by $N\approx\log\frac T{2\pi}$ at height $T$ and $N\approx\log c_f$ for a family & \raggedright Section~\ref{sec:ch6weyl}; Section~\ref{sec:ch6ks}\tabularnewline[1pt]
\raggedright $W(G)$ & \raggedright One level density of the symmetry type $G\in\{\mathrm U,\USp,\mathrm{SO(even)},\mathrm{SO(odd)},\mathrm O\}$: $\E\sum_j\phi(N\theta_j/2\pi)\to\int\phi\,W(G)$ & \raggedright Section~\ref{sec:ch6ks}\tabularnewline[1pt]
\raggedright $D(f;\phi)$, $c_f$ & \raggedright One level density statistic $\sum_{\gamma_f}\phi(\frac{\gamma_f}{2\pi}\log c_f)$ of the zeros $\frac12+i\gamma_f$ of $L(s,f)$, with analytic conductor $c_f$ & \raggedright Section~\ref{sec:ch6ks}; Theorem~\ref{thm:ch6ils}\tabularnewline[1pt]
\raggedright $\mathrm{Kl}(a,b;p)$, $K(t)$ & \raggedright Kloosterman sum $\sum_{x\in\mathbb F_p^\times}e((ax+b\bar x)/p)$, with $\mathrm{Kl}(a;p)/\sqrt p=2\cos\theta_{p,a}\in[-2,2]$; random Fourier series $K(t)=\sum_hc_h(t)\mathrm{ST}_h$ over independent $\mathrm{ST}_h$ with law $\ST$ & \raggedright Theorem~\ref{thm:06bweil}; Definition~\ref{def:06bK}\tabularnewline[1pt]
\raggedright $\mu_{\mathrm{CL}}$, $\eta_n(p)$ & \raggedright Cohen Lenstra measure $\mu_{\mathrm{CL}}(G)=\eta_\infty(p)/|\operatorname{Aut}G|$ on finite abelian $p$ groups, with $\eta_n(p)=\prod_{k=1}^n(1-p^{-k})$ & \raggedright Definition~\ref{def:06ccl}\tabularnewline[1pt]
\raggedright $U_p$, $X_p$ & \raggedright Independent uniform phases on the unit circle (the Steinhaus model) in the random Euler product $\prod_p(1-U_pp^{-s})^{-1}$; written $X_p$ in Chapter~\ref{ch:09} & \raggedright Theorem~\ref{thm:ch09bj}; Proposition~\ref{prop:09bcov}\tabularnewline[1pt]
\midrule
\multicolumn{3}{l}{\textit{Graphs and walks}}\\*
\raggedright $q$, $\T_q$ & \raggedright Prime (or integer $\ge2$) $q$, and the Bruhat Tits tree of $\PGL_2(\Q_q)$: the $(q+1)$ regular tree whose vertices are homothety classes of lattices in $\Q_q^2$ & \raggedright Chapter~\ref{ch:07}; Theorem~\ref{thm:07tree}\tabularnewline[1pt]
\raggedright $A$, $A_n$, $X$, $B_n$ & \raggedright On $\T_q$: adjacency operator $A=T_q$, sphere operators $A_nf(v)=\sum_{d(v,w)=n}f(w)$, $X=q^{-1/2}A$ and $B_n=q^{-n/2}A_n=U_n(X/2)-q^{-1}U_{n-2}(X/2)$ for $n\ge2$ & \raggedright Theorem~\ref{thm:07tree}\tabularnewline[1pt]
\raggedright $\mu_q$ & \raggedright Kesten McKay law, the spectral measure of $X$ at a vertex of $\T_q$, with $\mu_q\Rightarrow\ST$ as $q\to\infty$; written $\mu_\ell$, $\mu_p$ elsewhere & \raggedright Theorem~\ref{thm:07tree}; Chapters~\ref{ch:08} and \ref{ch:11b}\tabularnewline[1pt]
\raggedright $G_N(\ell)$, $B(\ell)$, $N$ & \raggedright Supersingular $\ell$ isogeny graph in characteristic $N$, a prime $N\equiv1\bmod12$ (the level of $S_2(\Gamma_0(N))$), $(\ell+1)$ regular on $(N-1)/12$ vertices, with adjacency matrix the Brandt matrix $B(\ell)$; $G_N(q)$ with $\ell=q$ in Chapter~\ref{ch:07}; in Chapter~\ref{ch:hit}, for any prime $N$, $P=B(\ell)/(\ell+1)$ is the isogeny walk, reversible for $\pi_i\propto1/w_i$ with $w_i=|\operatorname{Aut}(E_i)|/2$ & \raggedright Theorem~\ref{thm:08brandt}; Section~\ref{sec:hitisogeny}\tabularnewline[1pt]
\raggedright $h$ & \raggedright Number of vertices $(N-1)/12$ of $G_N(\ell)$, and for any prime $N$ the number $g+1$ of supersingular curves, $g$ the genus of $X_0(N)$ (Chapter~\ref{ch:hit}); also the class number $h(D)$ of Chapters~\ref{ch:06c} and \ref{ch:11b} & \raggedright Chapter~\ref{ch:08}; Theorem~\ref{thm:08kemeny}\tabularnewline[1pt]
\raggedright $P$, $\pi$ & \raggedright Transition matrix and stationary law of a reversible Markov chain; $P=A/d$ and $\pi$ uniform for a $d$ regular graph with adjacency matrix $A$; in Chapters~\ref{ch:lap} and \ref{ch:pd}, $P$ is a Dirichlet process (see $\Gamma_\mu$) & \raggedright Appendix~\ref{app:A}, Markov chains\tabularnewline[1pt]
\raggedright $\mu_j$, $\phi_j$ & \raggedright Eigenvalues $1=\mu_1>\mu_2\ge\dots\ge\mu_n\ge-1$ of $P$, with eigenvectors $\phi_j$ orthonormal in $L^2(\pi)$; $\mu_j=\lambda_j/d$ for adjacency eigenvalues $\lambda_j$. In Chapter~\ref{ch:10c} they are $p_j$, and $\mu_j=1-r_j$ are Dirichlet rates & \raggedright Appendix~\ref{app:A}, Markov chains; Chapter~\ref{ch:10c}\tabularnewline[1pt]
\raggedright $\tau_y$, $h(x,y)$ & \raggedright Hitting time $\tau_y=\min\{t\ge0:X_t=y\}$ and mean hitting time $h(x,y)=\E_x\tau_y$ & \raggedright Appendix~\ref{app:A}, Markov chains; Chapter~\ref{ch:08}\tabularnewline[1pt]
\raggedright $K$, $K_N$, $Z$ & \raggedright $K=\sum_y\pi(y)h(x,y)=\sum_{j\ge2}(1-\mu_j)^{-1}$, the Kemeny constant; fundamental matrix $Z=\sum_{t\ge0}(P^t-\Pi)$; $K_N$ for $G_N(\ell)$ with $K_N/h\to\ell/(\ell-1)$; $K_N(\ell)=\sum_f(\ell+1)/(\ell+1-a_\ell(f))$ for the isogeny walk at any prime $N$ & \raggedright Theorem~\ref{thm:A:kemeny}; Theorem~\ref{thm:08kemeny}; Theorem~\ref{thm:hitkemeny}\tabularnewline[1pt]
\raggedright $H_N(\ell)$ & \raggedright Hecke GGC with Thorin measure $\sum_f\delta_{\ell+1-a_\ell(f)}$ over the normalised eigenforms of $S_2(\Gamma_0(N))$, a sum of independent exponential variables with mean $K_N(\ell)/(\ell+1)$ & \raggedright Definition~\ref{def:flthecke}; Theorem~\ref{thm:fltempty}\tabularnewline[1pt]
\raggedright $\tau(G)$, $J_0(N)$, $C_N$ & \raggedright Number of spanning trees of a graph $G$, the order of its critical group; the Jacobian of $X_0(N)$, and for prime $N$ its cuspidal subgroup $C_N=J_0(N)(\Q)_{\mathrm{tors}}$, cyclic of order the numerator of $(N-1)/12$ & \raggedright Theorem~\ref{thm:hittrees}; Corollary~\ref{cor:hittorsion}\tabularnewline[1pt]
\raggedright $\rho$, $\rho_q$, $\rho_\ell$ & \raggedright Kesten rate $2\sqrt q/(q+1)$, the spectral radius of simple random walk on the $(q+1)$ regular tree; $\rho_\ell$ for $G_N(\ell)$ & \raggedright Theorem~\ref{thm:07tree}; Chapter~\ref{ch:08}\tabularnewline[1pt]
\raggedright $\lambda_*$, $\mu_*$ & \raggedright $\lambda_*=\max\{|\lambda_j|:|\lambda_j|\ne q+1\}$ for a $(q+1)$ regular graph; $\mu_*=\max(\mu_2,|\mu_n|)=\lambda_*/(q+1)$ if it is not bipartite & \raggedright Appendix~\ref{app:A}, Markov chains; Proposition~\ref{prop:08three}\tabularnewline[1pt]
\raggedright $\Lambda(\theta)$, $\Lambda^*(y)$ & \raggedright Limiting cumulant generating function $\Lambda(\theta)=\lim_nn^{-1}\log\E e^{n\theta Y_n}$ and rate function $\Lambda^*(y)=\sup_\theta(\theta y-\Lambda(\theta))$ & \raggedright Appendix~\ref{app:A}, G\"artner Ellis; Theorem~\ref{thm:10cldp}\tabularnewline[1pt]
\midrule
\multicolumn{3}{l}{\textit{Infinitely divisible laws, Thorin measures and GGCs}}\\*
\raggedright $(a,\beta^2,\nu)$ & \raggedright L\'evy Khintchine triple of an infinitely divisible law: drift $a$, Gaussian variance $\beta^2$ and L\'evy measure $\nu$ on $\R\setminus\{0\}$; in Lemma~\ref{lem:10dsin2} and Chapter~\ref{ch:lap}, $\nu$ and $\nu_\alpha$ are $x$ times the L\'evy measure of $\log f$ at the basepoint & \raggedright Theorem~\ref{thm:ch09lk}; Lemma~\ref{lem:10dsin2}\tabularnewline[1pt]
\raggedright $\tau$, $a$ & \raggedright Thorin measure $\tau$ and drift $a\ge0$ of a generalized gamma convolution (GGC), with $\E e^{-sX}=\exp(-as-\int\log(1+s/t)\,\tau(dt))$; its densities $U$, $U_\epsilon$, $u_\Gamma$ and $V_a$ (Chapter~\ref{ch:10d}; $T_\alpha=\frac12V_{\alpha-1/2}$ in Chapter~\ref{ch:beta}), the signed density $u_\varepsilon$ (Chapter~\ref{ch:scholze}) and $U_\alpha$ (Chapter~\ref{ch:flt}, the density $U_{\mathrm{ar}}$ of Chapter~\ref{ch:10d} for the Rankin Selberg product $D$ of $E$; in Chapter~\ref{ch:lap}, $U_{\mathrm{ar}}$ of any Euler product at the basepoint $\alpha$, signed when its coefficients are) and local combs $\tau_{\theta,L}$ (Chapter~\ref{ch:10c}); $\tau_Q$, the Thorin measure of a Frobenius functional (Theorem~\ref{thm:recipggcfun}); Tamagawa number $\tau(H)$ & \raggedright Definition~\ref{def:ch10ggc}; Equation \eqref{eq:10bggc}\tabularnewline[1pt]
\raggedright $k$, $k_\pi$, $\mathcal W_\pi$, $\mathcal W_\rho$, $\mathcal W_E$ & \raggedright L\'evy function: a GGC has L\'evy density $k(x)/x$ with $k(x)=\int e^{-tx}\tau(dt)$, and $k_\pi(x)=\sum_ke^{-t_kx}$ for $\pi$, written $\mathcal W_\pi$ (the heat trace) in Chapter~\ref{ch:10d}, where for self dual $\pi$ it is formed from $\Xi^{\mathrm{sd}}_\pi$; $\mathcal W_\rho$ is the Galois heat trace, defined from Frobenius traces, the conductor and the gamma factor, and $\mathcal W_E$ that of an elliptic curve (Chapter~\ref{ch:flt}) & \raggedright Equation \eqref{eq:10bkS}; Theorem~\ref{thm:ch10ggc}; Definition~\ref{def:recipheat}\tabularnewline[1pt]
\raggedright $t_k$, $S_n$ & \raggedright Thorin atoms $t_k=\gamma_k^2$, where $\pm i\gamma_k$ are the zeros of $\Xi_\pi$, and power sums $S_n=\sum_kt_k^{-n}$ & \raggedright Equations \eqref{eq:10batoms} and \eqref{eq:10bkS}\tabularnewline[1pt]
\raggedright $\xi$, $\Xi$, $\Phi$ & \raggedright Riemann's $\xi(s)=\frac12s(s-1)\pi^{-s/2}\Gamma(\frac s2)\zeta(s)$ and $\Xi(t)=\xi(\frac12+it)=\int\Phi(u)e^{itu}\,du$, with P\'olya's kernel $\Phi$; an unsubscripted $\Xi$ is always Riemann's & \raggedright Chapter~\ref{ch:10}; Equation \eqref{eq:ch10polya}\tabularnewline[1pt]
\raggedright $\Xi_\zeta$, $\Xi_\pi$, $\Xi^{\mathrm{sd}}_\pi$, $F_\pi$ & \raggedright $\Xi_\zeta(z)=\xi(\frac12+z)=\Xi(-iz)$; $\Xi_\pi(z)=\Lambda(\frac12+z,\pi)\Lambda(\frac12+z,\tilde\pi)/z^{2r}$, $r=\operatorname{ord}_{s=1/2}\Lambda(s,\pi)$; $\Xi^{\mathrm{sd}}_\pi=\Lambda(\frac12+z,\pi)/z^r$ if $\pi\cong\tilde\pi$, so $\Xi_\pi=(\Xi^{\mathrm{sd}}_\pi)^2$; $F_\pi(s)=\Xi_\pi(\sqrt s)/\Xi_\pi(0)=\prod_k(1+s/t_k)$; in Chapter~\ref{ch:scholze} the clock $G_\pi(s)=\Xi^{\mathrm{sd}}_\pi(\sqrt s)/\Xi^{\mathrm{sd}}_\pi(0)$ of a self dual $\pi$, so $G_\pi^2=F_\pi$ ($G_\zeta=F_\zeta$) & \raggedright Definition~\ref{def:10bxi}\tabularnewline[1pt]
\raggedright $\Phi_\pi$, $W_\pi$ & \raggedright P\'olya kernel with $\int\Phi_\pi(u)e^{zu}\,du=\Xi_\pi(z)/\Xi_\pi(0)$, and P\'olya variable $W_\pi$ with density $\Phi_\pi\ge0$; $W$ of Chapter~\ref{ch:10} uses $\Xi^{\mathrm{sd}}_\Delta$ and has density $\Delta(ie^u)e^{6u}/\Lambda(6,\Delta)$ (with $\Xi_\Delta$ it would be a sum of two copies); the kernels $\Phi_f$ and $\Phi_\chi$ of Chapter~\ref{ch:recip} and $\Phi_E$ of Chapter~\ref{ch:flt}, with P\'olya variable $W_E$, are not normalised: $\int\Phi_f(u)e^{tu}\,du=\Lambda(\frac k2+t)$; nor are $\Phi_\beta$ and $\Phi_{12}$ (Chapters~\ref{ch:beta} and \ref{ch:scholze}), $\Phi_K$ (Chapter~\ref{ch:squarefree}) and the kernel $\Phi_\rho=2\Phi_f$ of Chapter~\ref{ch:lap}, computed by Bromwich inversion, with $\int\Phi_\rho(u)e^{tu}\,du=\Lambda(\frac12+t,\rho)$, the factor $2$ coming from $\Gamma_\C$ & \raggedright Lemma~\ref{lem:10bkernel}; Theorem~\ref{thm:ch10kernel}; Equation \eqref{eq:recipmellin}; Proposition~\ref{prop:lappolya}\tabularnewline[1pt]
\raggedright $X$, $Y$, $Z_\pi$, $T$ & \raggedright The GGC $\sum_kE_k/t_k$ with drift $0$ and Thorin measure $\sum_k\delta_{t_k}$, named $X$ in Chapters~\ref{ch:10} and \ref{ch:10b}, $Y$ in Chapters~\ref{ch:09b} and \ref{ch:squarefree}, $Z_\pi$ in Chapter~\ref{ch:10d} and $T$ in Chapter~\ref{ch:hit} ($T_E$ in Chapter~\ref{ch:flt}, under GRH for $L(s,E)$), where it is the explosion time of a pure birth process with rates $t_k$; in Chapter~\ref{ch:squarefree}, $T=W_\zeta/2$ is half the P\'olya variable of $\zeta$, with moment generating function $M_T$ & \raggedright Theorem~\ref{thm:10babstract}; Theorem~\ref{thm:10dvdp}; Theorem~\ref{thm:hitexplosion}\tabularnewline[1pt]
\raggedright $E_k$, $L_k$ & \raggedright Independent standard exponential variables $E_k$ and Laplace variables $L_k$ (density $\frac12e^{-|x|}$); $E_k$ is also the Eisenstein series, $E_2$ quasimodular in Chapter~\ref{ch:11} & \raggedright Theorem~\ref{thm:ch10ggc}; Chapter~\ref{ch:10d}; Section~\ref{sec:02forms}\tabularnewline[1pt]
\raggedright $V$, $V_\pi$ & \raggedright Van Dantzig operator $VF(t)=1/F(it)$ on characteristic functions, and the series $V_\pi=\sum_kL_k/\gamma_k$; in Chapter~\ref{ch:beta}, $V=\operatorname{Var}X_\beta=2S_1$ for the P\'olya variable $X_\beta$ of $L(s,\chi_{-4})$ & \raggedright Definition~\ref{def:10dvdp}; Theorem~\ref{thm:10dvdp}\tabularnewline[1pt]
\raggedright $(X,H)$ & \raggedright Wald couple: $H\ge0$ and $\E e^{sX}\,\E e^{-s^2H}=1$ for small $s\ge0$; under the Riemann hypothesis $(W_\pi,Z_\pi)$ is one & \raggedright Definition~\ref{def:10dwald}\tabularnewline[1pt]
\raggedright $H$, $H_a$, $C$ & \raggedright The GGC $H$ of the sine squared lemma, $\E e^{-sH}=e^{b\sqrt s}f(\alpha)/f(\alpha+\sqrt s)$ with $b=(f'/f)(\alpha)$, whose tail Chapter~\ref{ch:lap} computes; the GGC $H_a$ with $\E e^{-sH_a}=\Xi_\pi(a)/\Xi_\pi(a+\sqrt s)$ and Thorin density $V_a$, formed in Chapter~\ref{ch:lap} from $\Xi^{\mathrm{sd}}_\pi$, where $V_a$ is, in the variable $r=\sqrt z$, a superposition over the zeros of folded Cauchy laws, $C$ being a standard Cauchy variable & \raggedright Lemma~\ref{lem:10dsin2}; Theorem~\ref{thm:10dshift}; Theorems~\ref{thm:lapcauchy} and \ref{thm:laptail}\tabularnewline[1pt]
\raggedright $\Theta_\sigma$, $P_\sigma$ & \raggedright Parabolic Thorin condition: every zero of $\Lambda(s,\pi)$ has $|\Re s-\frac12|\le\sigma$, equivalently every atom lies in $-P_\sigma$; $\Theta_0$ is the Riemann hypothesis, $\Theta_{1/2}$ is unconditional & \raggedright Definition~\ref{def:10btheta}\tabularnewline[1pt]
\raggedright $H_{\pi,\lambda}$, $\lambda_\pi$, $\Lambda$ & \raggedright Gaussian tilt $H_{\pi,\lambda}(t)=\int e^{\lambda u^2}\Phi_\pi(u)e^{itu}\,du$, and the de Bruijn Newman constant $\lambda_\pi\le\frac18$ in the book's units; the literature constant is $\Lambda=4\lambda_0$ for $\zeta$, with $\Lambda\le\frac12$ & \raggedright Theorem~\ref{thm:10btilting}; Section~\ref{sec:10bparabola}; Chapter~\ref{ch:10}\tabularnewline[1pt]
\raggedright $N_\sigma$, $e_p$, $N^{(\vartheta)}_\sigma$ & \raggedright Zeta distributed random integer, $\Prob(N_\sigma=n)=n^{-\sigma}/\zeta(\sigma)$, with independent geometric exponents $e_p$ (written $N$ and $G_p$ in Chapter~\ref{ch:09}); $\log N_\sigma$ is the zeta subordinator at time $1$, with L\'evy measure $\Pi_\sigma$; $N^{(\vartheta)}_\sigma$, whose logarithm is the zeta subordinator at time $\vartheta$, has the law $d_\vartheta(n)n^{-\sigma}/\zeta(\sigma)^\vartheta$, with $d_\vartheta$ the coefficients of $\zeta^\vartheta$ & \raggedright Theorem~\ref{thm:ch09khinchine}; Proposition~\ref{prop:10dsublaw}; Lemma~\ref{lem:recipzeta}; Theorem~\ref{thm:pdtheta}\tabularnewline[1pt]
\raggedright $\Gamma_\mu$, $\Xi_\sigma$, $\mathcal M_\sigma$, $\mathcal P_\sigma$ & \raggedright Gamma random measure with shape measure $\mu$; the Frobenius labelled random measure of $(\sigma-1)\log N_\sigma$ (a subscripted $\Xi$ that is not a completion) and the marked point process of its prime factors; normalised, $\mathcal P_\sigma=\Xi_\sigma/\Xi_\sigma(Y)$ (and $\mu$ when $\Xi_\sigma=0$), which converges in law to the Dirichlet process $P$ with base measure $\mu$ & \raggedright Theorem~\ref{thm:recipgammaproc}; Theorem~\ref{thm:pddirichlet}\tabularnewline[1pt]
\raggedright $Q_{E,p}$, $c(E,p)$, $c_p$ & \raggedright Eisenstein discrepancy $Q_{E,p}(\ell)=\mathsf 1\{p\nmid|E(\mathbb F_\ell)|\}$ and Eisenstein shape $c(E,p)=\mu_{G_{E,p}}\{g:\det(1-g)\ne0\}$, the local density of Koblitz's conjecture and the shape of the Gamma limit of $\Xi_\sigma(Q_{E,p})$; $c_p=1-(p^2-2)/((p^2-1)(p-1))$ when $G_{E,p}=\GL_2(\mathbb F_p)$ & \raggedright Definition~\ref{def:flteis}; Proposition~\ref{prop:flteisgamma}\tabularnewline[1pt]
\raggedright $\mathrm{PD}(\vartheta)$, $V_k$, $W_k$ & \raggedright Poisson Dirichlet law with parameter $\vartheta$: the law of the ranked jumps of a standard Gamma subordinator on $[0,\vartheta]$ divided by their sum, written $V_1\ge V_2\ge\dots$; for $\vartheta=1$ the stick breaking sequence $W_k=(1-U_1)\cdots(1-U_{k-1})U_k$, with independent uniform $U_i$, is its size biased order (the GEM law); not the operator $V$, the densities $V_a$ or the P\'olya variable $W$ & \raggedright Section~\ref{sec:pdlaw}; Theorem~\ref{thm:pdtheta}\tabularnewline[1pt]
\midrule
\multicolumn{3}{l}{\textit{Extreme values and eta products}}\\*
\raggedright $\eta(\tau)$ & \raggedright Dedekind eta function $q^{1/24}\prod_{n\ge1}(1-q^n)$, and eta products $f=\prod_m\eta(m\tau)^{r_m}$ with $\sum_mm\,r_m=24$ & \raggedright Chapter~\ref{ch:11}; Theorem~\ref{thm:ch11general}\tabularnewline[1pt]
\raggedright $M$, $P(y)$ & \raggedright $M=\max_{n\ge1}E_n/(2\pi n)$ for standard exponentials $E_n$, so $\Prob(M\le y)=P(y)=\prod_{n\ge1}(1-e^{-2\pi ny})$; $M$ is also a clipping level (Chapter~\ref{ch:12}), $M(x;u_0,\dots,u_k)$ is the Curry Schoenberg B spline with knots $u_0,\dots,u_k$ (Chapter~\ref{ch:lap}), and $M(t)=\sum_{d\le t}\mu(d)$ is the Mertens function in Chapter~\ref{ch:squarefree} & \raggedright Proposition~\ref{prop:ch11M}; Chapter~\ref{ch:10c}; Section~\ref{sec:lapspline}\tabularnewline[1pt]
\raggedright $M^*$, $M_f$, $F$ & \raggedright Maximum $M^*$ of $24$ independent copies of $M$, with $F(y)=\Prob(M^*\le y)=P(y)^{24}$ and $\Delta(iy)=e^{-2\pi y}F(y)$; $M_f$ is the analogue for an eta product $f$ & \raggedright Theorem~\ref{thm:ch11main}; Theorem~\ref{thm:ch11general}\tabularnewline[1pt]
\raggedright $Q(s,c)$, $G_s$ & \raggedright Regularised incomplete gamma function $\Gamma(s,c)/\Gamma(s)$, and a gamma variable $G_s$ of shape $s$ independent of $M^*$: $L(s,\Delta)=\E\,Q(s,2\pi M^*)=\Prob(G_s\ge2\pi M^*)$ for $s>0$ in the classical normalisation & \raggedright Theorem~\ref{thm:ch11main}\tabularnewline[1pt]
\midrule
\multicolumn{3}{l}{\textit{Statistics and inference}}\\*
\raggedright $D_n$, $F_n$, $\mathcal K$ & \raggedright Kolmogorov discrepancy $D_n=\sup_t|F_n(t)-F(t)|$ of the empirical distribution $F_n$ of $x_1,\dots,x_n$; $\sqrt n\,D_n\Rightarrow\mathcal K$ (Kolmogorov law) for independent draws ($K$ in Chapters~\ref{ch:05}, \ref{ch:11}) & \raggedright Appendix~\ref{app:A}, equidistribution; Section~\ref{sec:05disc}\tabularnewline[1pt]
\raggedright $\mathrm{KL}(\mu\|\nu)$ & \raggedright Kullback Leibler divergence $\int\log(d\mu/d\nu)\,d\mu$, equal to $+\infty$ unless $\mu\ll\nu$ & \raggedright Appendix~\ref{app:A}, Bayes factors\tabularnewline[1pt]
\raggedright $B_n$, $B_n^M$, $h$, $h_M$ & \raggedright Bayes factor $B_n=\prod_{i\le n}f_1(X_i)/f_0(X_i)$ of $f_1$ against $f_0$, not the sphere operators $B_n$; with the log likelihood ratio $h=\log(f_1/f_0)$ and its clipping $h_M=\max(-M,\min(M,h))$, $\log B_n^M=\sum_{k\le n}h_M(x_k)$ & \raggedright Appendix~\ref{app:A}, Bayes factors; Theorem~\ref{thm:12consist}\tabularnewline[1pt]
\raggedright $M_{\mathrm{ST}}$, $M_{\mathrm{CM}}$, $\varepsilon$ & \raggedright Models $M_{\mathrm{ST}}=(1-\varepsilon)\ST+\varepsilon\delta_0$ and $M_{\mathrm{CM}}=\mu_{\mathrm{CM}}$ for the traces $x_p$, with a small atom $\varepsilon$ at $0$ & \raggedright Corollary~\ref{thm:12elliptic}\tabularnewline[1pt]
\bottomrule
\end{longtable}}
